% SUBSTITUTE for the "Lachaise Assignment" template (originally Vel /
% mursisoy, CC-BY-NC-SA — non-commercial, so it cannot be redistributed
% here).
%
% This starter uses the upstream `homework` document class (CTAN, part
% of Jinwen XU's minimart/einfart class family, LPPL-1.3c) — built for
% exactly this purpose: problem sets with numbered problems, solutions,
% proofs and QED marks. See README.md for the full disclosure.
%
% homework is not bundled: it ships in TeX Live via its own CTAN package.
\documentclass[11pt,
  title in boldface,
  theorem in new line,
  simple names,
  colored solution,
]{homework}

\title{Problem Set 4: Sequences and Series}
\author{Prof.\ Ana Petrova}
\date{\TheDate{2026-02-10}, Numerical Methods}

\begin{document}

\begin{problem}
Determine whether the sequence $a_n = \dfrac{n^2 + 3n}{2n^2 - 1}$
converges, and if so, find its limit.
\end{problem}

\begin{solution}
Divide numerator and denominator by $n^2$:
\[
  a_n = \frac{1 + 3/n}{2 - 1/n^2} \longrightarrow \frac{1}{2}
  \quad \text{as } n \to \infty.
\]
The sequence converges to $\tfrac{1}{2}$.
\end{solution}

\begin{problem}[Ratio test]
Use the ratio test to determine convergence of
$\displaystyle\sum_{n=1}^{\infty} \frac{n!}{n^n}$.
\end{problem}

\begin{solution}[Proof]
Let $a_n = n!/n^n$. Then
\[
  \frac{a_{n+1}}{a_n}
  = \frac{(n+1)!}{(n+1)^{n+1}} \cdot \frac{n^n}{n!}
  = \left(\frac{n}{n+1}\right)^n
  = \left(1 + \frac{1}{n}\right)^{-n}
  \longrightarrow e^{-1} < 1.
\]
By the ratio test, the series converges.
\qedhere
\end{solution}

\begin{problem}[Several parts]
Consider the sequence defined by $x_1 = 1$ and $x_{n+1} = \sqrt{2 + x_n}$.
\begin{enumerate}
  \item Show that $x_n < 2$ for all $n$.

  \begin{solution}
  By induction: $x_1 = 1 < 2$. If $x_n < 2$, then
  $x_{n+1} = \sqrt{2 + x_n} < \sqrt{4} = 2$.
  \end{solution}

  \item Show that $(x_n)$ is increasing.

  \begin{solution}
  Since $x_n < 2$, we have $x_n^2 < 2x_n + \text{(check sign)}$, which
  after simplification gives $x_{n+1} > x_n$ for all $n$.
  \end{solution}

  \item Conclude that $(x_n)$ converges, and find its limit.

  \begin{solution}
  A bounded, monotonically increasing sequence converges. Taking
  limits in the recursion, $L = \sqrt{2 + L}$ gives $L^2 - L - 2 = 0$,
  so $L = 2$.
  \end{solution}
\end{enumerate}
\end{problem}

\begin{fact}
A convergent sequence is bounded.
\proofless
\end{fact}

\DNF<the divergence test for the harmonic series>

\end{document}
